\section{Machines from blocks}
\label{sec:blocks}

From here to \cref{sec:demos} the paper is a plan. Nothing in \cref{sec:blocks,sec:arch,sec:runtime} is built;
\cref{sec:demos} registers, for each piece, what we expect to measure and what result would make us abandon it.
We write the plan as a paper because the design decisions are the research: once the calculus below is fixed,
building it is engineering, and if the calculus is wrong no amount of training will rescue it.

\subsection{Blocks and ports}
\label{sec:ports}

A \emph{block} is a spec, its network, its owned fields and three kinds of port. It \emph{owns} its state
fields: no other block writes them. It \emph{accepts} actions, which are its input ports; some arrive from a
user, some from other blocks. It \emph{emits} events, declared in the spec as a new clause on a transition
(\texttt{emits: \{ wrap: "old.d == 9" \}}, a predicate over the old and new state that the runtime evaluates
on the step the block just took), which are its output ports. And it may \emph{read} fields of other blocks, named in the spec as
\texttt{reads:}; a read field is part of the block's local state for every purpose except ownership. A block's
\emph{cells} are the pairs (local state, action) over its own fields, its read fields, and its accepted actions.

A block is \emph{exact} when its network's step equals its oracle's step on every cell and breaks no invariant
of the block. That is \cref{tab:round2}'s result, restated per block. Its \emph{contract} is its invariant and
its \texttt{ensures}; with valid candidates (\cref{sec:valid}) the invariant half holds by construction at the
network's output, whatever the weights, and the \texttt{ensures} half is what training is held to.

\subsection{Three couplings and what each costs}
\label{sec:couplings}

How two blocks are joined decides what each one's network has to learn.

\paragraph{Independent.} The blocks share nothing. Their energies add, so the machine's argmin splits into one
argmin per block, and the machine's cost is the sum of the blocks' costs. The scheduler of \cref{tab:round2} is
three copies of a twelve-state meeting block that happen to be trained as one $1{,}728$-state network; as a
machine it is one network of twelve states, used three times.

\paragraph{Event-wired.} A step of block $A$ emits an event that the wiring delivers as an action to block $B$.
$B$'s cells do not mention $A$; $B$ learns \texttt{inc}, not \emph{why} it was incremented. The machine's state
space is the product of the blocks' spaces, but nothing ever enumerates it: a six-digit odometer has $10^6$
states and six copies of a ten-state block. Cost is again the sum. This is the coupling hardware is built from,
a register and a carry line, and it is the one we expect to carry most of the weight.

\paragraph{Read-wired.} $B$'s rule reads $A$'s fields. $B$'s cells are now its own states times the read
fields' domains times its actions, so the cost grows with the coupling: Sokoban's box reads the player's cell
and the cell beyond, which is still small; kanban's every card reads every other card, so the board does not
split at all and its one network is the right design. The rule for a designer is the rule for a hardware
designer: keep read ports few and narrow, and move coupling into events wherever a read can be replaced by a
message. We expect that most interfaces people actually use are closer to the odometer than to kanban.

\subsection{Exactness composes}
\label{sec:theorem}

\begin{definition}[Machine]
A machine is a finite set of blocks $B_1,\dots,B_n$ with specs $S_i$, a wiring $W$ that maps each emitted
event $(i,e)$ to a finite list of deliveries $(j,a_j)$ with $a_j$ an action of $S_j$ whose parameters are
computed from $e$'s payload, and a read relation giving each block the fields it reads. A \emph{tick} applies
a finite ordered list of external actions, each to its block, and after each one processes the resulting event
queue in a fixed order, each delivery stepping its target block in the machine state current at that moment,
until the queue is empty. Emitted events are computed by the runtime from the spec's \texttt{emits} clauses on
the block's old and new states; they are derived, not learned. The
\emph{composed spec} $S_M$ is this procedure with every block stepped by its oracle; the \emph{composed
network} $N_M$ is the same procedure with every block stepped by its network.
\end{definition}

\begin{theorem}[Exactness composes]
\label{thm:compose}
If every block is exact on its cells and every tick's event cascade terminates, then $N_M(s,a) = S_M(s,a)$ for
every reachable machine state $s$ and external action $a$, and no invariant of any block is broken after any
tick.
\end{theorem}

\begin{proof}[Proof sketch]
By induction on the cascade, one external action at a time. The action steps block $i$ from a local state
(its own fields and its read fields) that is identical under $N_M$ and $S_M$ before it; exactness gives the
same successor, and the same successor gives the same emitted events, because events are evaluated from the
spec on the states and not produced by the network. Each delivery then steps its target in a machine state that, by the inductive
hypothesis, is identical under both, with the same action and parameters, so again the same successor and the
same events; the queues stay equal, and termination makes the induction finite. Invariants: every block step is
a transition of that block's oracle, which preserves its invariant, and with valid candidates the network
cannot leave the invariant's set even where it is wrong.
\end{proof}

\begin{remark}
The theorem is about equality with the composed spec, not about the composed spec being right. Wiring is code
again, and a wiring bug is a bug; the composed oracle catches it only in the sense that a test suite catches
bugs, by running the machine against its own invariants and against properties that span blocks (\texttt{the
odometer's value equals the number of increments}), exhaustively when the reachable set is small and by
sampling when it is not. Liveness and other temporal properties are outside the theorem and belong to the
bounded-LTL checker that Part I's roadmap already lists.
\end{remark}

\begin{remark}
Termination is a property of the wiring graph. An acyclic event graph terminates trivially; a cycle (two
blocks that notify each other) needs a bound per tick, and we propose that the spec declares one and the
checker enforces it, as synchronous languages require causality to be resolved at compile time~\citep{berry1992esterel,benveniste2003synchronous}.
\end{remark}

\subsection{The composed spec}
\label{sec:composed}

The spec language gains a \texttt{machine} document that names block types, instances, wires and reads.
\Cref{fig:odometer-spec} writes the odometer. Instances of one block type share one network; the composed
oracle is derived from the block oracles and the wiring, like every oracle in \method, and \texttt{nerve check}
runs the machine against itself before anything is trained.

\begin{figure}[t]
\begin{minipage}[t]{0.47\textwidth}
\begin{lstlisting}
nerve: 1
name: digit
state:
  d: {type: int, range: [0, 9]}
actions:
  inc: {}
  reset: {}
transitions:
  inc:   { d: "(d + 1) % 10",
           emits: { wrap: "old.d == 9" } }
  reset: { d: "0" }
ensures:
  - { after: inc, holds: "d == (old.d + 1) % 10" }
\end{lstlisting}
\end{minipage}\hfill
\begin{minipage}[t]{0.47\textwidth}
\begin{lstlisting}
nerve: 1
machine: odometer
blocks:
  d0: digit
  d1: digit
  d2: digit
wires:
  - on: d0.wrap
    do: d1.inc
  - on: d1.wrap
    do: d2.inc
ensures:                 # spans blocks
  - { after: d0.inc,
      holds: "value == (old.value + 1) % 1000" }
derived:
  value: "d0.d + 10*d1.d + 100*d2.d"
\end{lstlisting}
\end{minipage}
\caption{\textbf{A machine is blocks plus wiring.} Left: a ten-state digit block that emits \texttt{wrap} when
it rolls over. Right: three instances and two carry wires make a thousand-state odometer whose network is the
digit's, used three times. The machine's own \texttt{ensures} is checked by the composed oracle, exhaustively
at this size.}
\label{fig:odometer-spec}
\end{figure}
